% =====================================================================
%  Strona tytułowa — MorphBPE: Morfologicznie informowany BPE dla polszczyzny
% =====================================================================
\begin{titlepage}
\centering
\vspace*{0.8cm}

\slayertitle
  {MorphBPE --- morfologicznie informowany Byte-Pair Encoding dla języka polskiego}%
  {Piotr Styła}%
  {\today}

\vspace{0.9cm}

% --- WYKRES HERO: MorphED — BPE vs MorphBPE przy vocab=500 i 1000 ---
\begin{figure}[h]
  \centering
  \begin{tikzpicture}
    \begin{axis}[
        ybar,
        bar width=16pt,
        width=0.82\linewidth, height=6.5cm,
        symbolic x coords={vocab=500, vocab=1000},
        xtick=data,
        ymin=1.5, ymax=2.7,
        ylabel={Morphological Edit Distance (niżej = lepiej)},
        nodes near coords,
        every node near coord/.style={font=\small},
        xmajorgrids=true, ymajorgrids=true,
        grid style={gray!20},
        legend style={at={(0.5,-0.25)}, anchor=north, legend columns=2, font=\small},
        enlarge x limits=0.5,
    ]
      \addplot[fill=slayerblue!60, draw=slayerblue] coordinates {
        (vocab=500, 2.503)
        (vocab=1000, 1.838)
      };
      \addlegendentry{BPE}
      \addplot[fill=slayergreen!60, draw=slayergreen] coordinates {
        (vocab=500, 2.414)
        (vocab=1000, 1.712)
      };
      \addlegendentry{MorphBPE suffix}
    \end{axis}
  \end{tikzpicture}
  \caption{\textbf{Morphological Edit Distance (MorphED) --- BPE vs MorphBPE.}
  Korpus: 10\,000 przemówień sejmowych (80/20 split, seed=42).
  Bootstrap CI 95\%: MorphBPE suffix vs BPE przy vocab=500:
  $2{,}414\,[2{,}408,\,2{,}420]$ vs $2{,}503\,[2{,}497,\,2{,}509]$
  ($-3{,}6\%$); przy vocab=1000:
  $1{,}712\,[1{,}707,\,1{,}717]$ vs $1{,}838\,[1{,}832,\,1{,}843]$
  ($-6{,}8\%$). Niżej = lepszy alignment morfologiczny.}
  \label{fig:hero}
\end{figure}

\vfill
{\small Piotr Styła \quad Polish Sejm Corpus \quad Slayer AI Lab \quad kompilowany w~CI}
\end{titlepage}
